\documentclass[12pt, oneside]{amsart}

% ── Fonts & encoding ──────────────────────────────────────────────────────────
\usepackage[T1]{fontenc}
\usepackage{mathpazo}
\usepackage{microtype}

% ── Mathematics ───────────────────────────────────────────────────────────────
\usepackage{amsmath, amssymb, amsthm, mathtools}
\usepackage{stmaryrd}     % \llbracket, \rrbracket

% ── Layout ────────────────────────────────────────────────────────────────────
\usepackage[margin=1.25in, top=1.4in, bottom=1.4in]{geometry}
\usepackage{setspace}
\onehalfspacing
\usepackage{titlesec}
\usepackage{fancyhdr}
\usepackage{enumitem}
\usepackage{hyperref}
\hypersetup{
  colorlinks=true,
  linkcolor=black,
  citecolor=black,
  urlcolor=black
}

% ── Theorem environments ───────────────────────────────────────────────────────
\theoremstyle{plain}
\newtheorem{theorem}{Theorem}[section]
\newtheorem{lemma}[theorem]{Lemma}
\newtheorem{proposition}[theorem]{Proposition}
\newtheorem{corollary}[theorem]{Corollary}

\theoremstyle{definition}
\newtheorem{definition}[theorem]{Definition}
\newtheorem{construction}[theorem]{Construction}
\newtheorem{example}[theorem]{Example}

\theoremstyle{remark}
\newtheorem{remark}[theorem]{Remark}
\newtheorem{notation}[theorem]{Notation}

% ── Custom operators and macros ───────────────────────────────────────────────
\DeclareMathOperator{\Vol}{Vol}
\DeclareMathOperator{\Res}{Res}
\DeclareMathOperator{\Int}{Int}
\DeclareMathOperator{\cl}{cl}
\DeclareMathOperator{\bd}{bd}
\DeclareMathOperator{\diam}{diam}
\newcommand{\R}{\mathbb{R}}
\newcommand{\N}{\mathbb{N}}
\newcommand{\Z}{\mathbb{Z}}
\newcommand{\C}{\mathbb{C}}
\newcommand{\Gr}{\mathrm{Gr}}
% \Phi already defined
\newcommand{\Om}{\Omega}
\newcommand{\om}{\omega}
\newcommand{\eps}{\varepsilon}
\newcommand{\dd}{\,\mathrm{d}}
\newcommand{\sfrac}[2]{{\textstyle\frac{#1}{#2}}}

% RSVP field operators
\newcommand{\Phiop}{\widehat{\Phi}}
\newcommand{\vop}{\widehat{v}}
\newcommand{\Sop}{\widehat{S}}
\newcommand{\Kop}{\mathcal{K}}

% KES operators
\newcommand{\XKS}{\mathcal{X}}
\newcommand{\OmKS}{\Omega}
\newcommand{\HKS}{\mathcal{H}}
\newcommand{\EKS}{\mathcal{E}}

% Spherepop
\newcommand{\Spop}{\mathbf{Sp}}
\newcommand{\hist}{\mathfrak{h}}
\newcommand{\Log}{\mathfrak{L}}

% Amplitude and geometry
\newcommand{\Amp}{\mathcal{A}}
\newcommand{\WF}{\Psi}
\newcommand{\Assoc}{\mathcal{K}_n}
\newcommand{\Cosmo}{\mathcal{C}_n}
\newcommand{\Water}{\mathcal{W}_n}
\newcommand{\PG}{\mathcal{P}}
\newcommand{\CF}{\Omega}   % canonical form

% ── Headers ───────────────────────────────────────────────────────────────────
\pagestyle{fancy}
\fancyhf{}
\renewcommand{\headrulewidth}{0.3pt}
\fancyhead[L]{\small\itshape Plenum, Event, and Positive Geometry}
\fancyhead[R]{\small\thepage}

% ── Title block ───────────────────────────────────────────────────────────────
\title{\textbf{Plenum, Event, and Positive Geometry}\\[0.5em]
\large Deriving Amplituhedra, Associahedra, Cosmohedra,\\
and Waterhedra from RSVP, KES, and Spherepop}

\author{Flyxion}
\date{Independent Researcher \\ \smallskip 2026}

% ══════════════════════════════════════════════════════════════════════════════
\begin{document}
% ══════════════════════════════════════════════════════════════════════════════

\maketitle
\thispagestyle{empty}

\begin{abstract}
We demonstrate that the family of positive geometries introduced by Arkani-Hamed and collaborators to reformulate scattering amplitudes and cosmological correlators without reference to spacetime---the associahedron, the amplituhedron, the cosmohedron, and the recently conjectured waterhedra---arise naturally and necessarily from three independently motivated theoretical frameworks: the Relativistic Scalar-Vector Plenum (RSVP), the Kinetic-Event Synthesis (KES), and the Spherepop event calculus. In RSVP, the scalar constraint-density field $\Phi$ defines a viability region in field configuration space whose boundary structure is shown to recover the associahedron as a special case of a constraint polytope. In KES, the possibility space $\OmKS_t$ and irreversible event history $\HKS_t$ jointly define a canonical differential form on kinematic space whose singularity locus coincides with the boundary stratification of positive geometries in general. In Spherepop, the append-only event log furnishes a history functor that extends the cosmohedron's face poset to an ordered categorical structure encoding the full arrow of cosmological time. We further conjecture, and provide initial evidence for, a waterhedron construction arising from the RSVP velocity field $v$ restricted to one-dimensional wave propagation, matching the structure recently observed by Arkani-Hamed and collaborators for deep-water wave amplitudes. The unifying principle across all four derivations is that \emph{combinatorics is the residue of irreversibility}: when spacetime and its metric are absent, what remains in any history-first ontology is precisely the positivity and boundary stratification that characterises these geometric objects. This places RSVP/KES/Spherepop in the role of a generative foundation for the positive geometry programme rather than a mere analogy to it.
\end{abstract}

\vspace{1em}
\tableofcontents

\newpage

% ══════════════════════════════════════════════════════════════════════════════
\section{Introduction}
% ══════════════════════════════════════════════════════════════════════════════

The positive geometry programme, as articulated most fully in the work of Arkani-Hamed and collaborators over the past fifteen years, pursues a radical reorganisation of fundamental physics. Its central wager is that locality and unitarity---long treated as axiomatic within quantum field theory---are not primitive laws but emergent consequences of more fundamental mathematical structures: combinatorial orderings, positivity conditions, and the canonical differential forms associated to geometric regions in kinematic space. When spacetime is stripped away and one operates directly in the abstract space of particle momenta, a ``sliver of structure'' remains, and it is this sliver that determines scattering amplitudes, cosmological correlators, and, apparently, the amplitudes of ocean waves.

The present paper argues that this sliver has a precise name and a rigorous generative source. It is the residue of irreversibility in the sense formalised by the RSVP, KES, and Spherepop frameworks developed independently by this author. The argument proceeds by explicit derivation rather than analogy: we show that starting from the field equations of RSVP, the ontological commitments of KES, and the categorical structure of the Spherepop event calculus, each of the major geometric objects of the positive geometry programme falls out as a theorem or a well-defined construction, and that Spherepop provides a natural extension that the positive geometry programme currently lacks.

\subsection{The Positive Geometry Programme}

Let us briefly characterise the main objects at stake. A positive geometry is a pair $(\mathcal{P}, \CF)$ where $\mathcal{P}$ is a region in some ambient space $V$ possessing boundaries of all codimensions, and $\CF(\mathcal{P})$ is a rational differential form, the \emph{canonical form}, uniquely characterised by the property that its singularities occur if and only if one approaches a boundary of $\mathcal{P}$, with the correct residue structure at each boundary component. The physical content is then the claim that scattering amplitudes are canonical forms of geometries living directly in kinematic space, with no reference to virtual particles, internal propagators, or spacetime locality.

The four principal examples treated here are:
\begin{itemize}[noitemsep]
  \item The \textbf{associahedron} $\Assoc$, which encodes tree-level scattering of ordered coloured scalars; its face poset is the poset of non-crossing chord diagrams on an $n$-gon.
  \item The \textbf{amplituhedron}, which encodes all-loop gluon scattering in $\mathcal{N}=4$ super-Yang-Mills; it lives in the Grassmannian $\Gr(k, n)$ and its positive region.
  \item The \textbf{cosmohedron} $\Cosmo$, announced in December 2024, which encodes the cosmological wave function; it is obtained by shaving the associahedron with additional inequalities arising from energy denominators.
  \item The \textbf{waterhedra} $\Water$, a family of geometric objects conjectured in early 2026 to govern deep-water wave amplitudes in one spatial dimension.
\end{itemize}

\subsection{The Generative Frameworks}

The three frameworks from which we derive these objects are sketched here; full technical treatments appear in the main sections.

The \textbf{Relativistic Scalar-Vector Plenum} (RSVP) posits that physical reality is constituted by three coupled fields on a spacetime manifold $M$: a scalar constraint-density field $\Phi: M \to \R$, a vector field $v: M \to TM$, and an entropy field $S: M \to \R$. The field equations couple these through a constitutive relation that makes the zero-set of $\Phi$ the boundary of the physically accessible region, and makes $S$ a Lyapunov function for the dynamics restricted to that region.

The \textbf{Kinetic-Event Synthesis} (KES) is a process ontology in which the fundamental state of a system at time $t$ is the triple $(\XKS_t, \OmKS_t, \HKS_t)$ where $\XKS_t$ is the current configuration, $\OmKS_t$ is the possibility space of available transitions, and $\HKS_t$ is the irreversible event history. The central claim is that physical observables are determined by the geometry of $\OmKS_t$ and the combinatorial structure of $\HKS_t$, not by $\XKS_t$ directly.

The \textbf{Spherepop event calculus} models identity and structure as append-only sequences of irreversible events. An entity is its history $\hist = (e_1, e_2, \ldots, e_k)$ together with a log functor $\Log$ mapping the history poset to a category of observable consequences. The defining commitments are that events are irreversible, histories are append-only, and the only admissible operations are those that extend rather than rewrite the log.

\subsection{The Unifying Thesis}

The thesis of this paper, stated without proof in the abstract and established section by section below, is the following. When one takes seriously the history-first ontology common to KES and Spherepop, and the boundary-first conception of physical possibility embodied in RSVP, the combinatorial residue that remains after spacetime structure is quotiented out is precisely the combinatorics of non-crossing partitions, nested tubings, and positive half-spaces that Arkani-Hamed's programme identifies as the foundation of physics. Locality and unitarity emerge from irreversibility and constraint not by analogy but by derivation. The cosmohedron's extra structure---the shaving of the associahedron by energy denominator inequalities---is identified as the imprint of the Spherepop log on the possibility space: histories that are causally ordered impose additional inequalities that carve out the cosmohedron from within the associahedron. And the waterhedra arise from the restriction of the RSVP velocity field to one spatial dimension, where the wave-like constitutive equation reduces to an integrable system whose amplitude is computed by a volume form on a convex polytope.

\subsection{Organisation of the Paper}

Section~\ref{sec:rsvp} develops the RSVP framework in sufficient technical detail and derives the associahedron from its constraint geometry. Section~\ref{sec:kes} develops KES and derives the general notion of a positive geometry with its canonical form. Section~\ref{sec:spherepop} develops the Spherepop calculus and derives the cosmohedron as the positive geometry of a history-enriched kinematic space. Section~\ref{sec:water} constructs the waterhedra from the RSVP velocity field in one dimension and establishes the correspondence with the structures observed for deep-water waves. Section~\ref{sec:unification} states the master theorem unifying all four derivations and discusses the implications for the positive geometry programme. Section~\ref{sec:discussion} addresses open questions, limitations of the present treatment, and directions for future work.

% ══════════════════════════════════════════════════════════════════════════════
\section{RSVP and the Associahedron}
\label{sec:rsvp}
% ══════════════════════════════════════════════════════════════════════════════

\subsection{The RSVP Field Equations}

We begin by fixing the RSVP field content precisely.

\begin{definition}[RSVP Triple]
\label{def:rsvp}
Let $M$ be a smooth $d$-dimensional Lorentzian manifold. An \emph{RSVP triple} on $M$ is a collection $(\Phi, v, S)$ where:
\begin{enumerate}[label=(\roman*)]
  \item $\Phi \in C^\infty(M, \R)$ is the \emph{scalar constraint-density field};
  \item $v \in \Gamma(TM)$ is the \emph{plenum velocity field};
  \item $S \in C^\infty(M, \R_{\geq 0})$ is the \emph{entropy field};
\end{enumerate}
subject to the \emph{RSVP field equations}:
\begin{align}
  \Box \Phi + \nabla_v \Phi &= -\lambda S \Phi, \label{eq:phi}\\
  \nabla_v v + \nabla \Phi &= -\kappa v, \label{eq:v}\\
  \partial_t S + \nabla \cdot (Sv) &= \sigma \|\nabla \Phi\|^2, \label{eq:S}
\end{align}
where $\lambda, \kappa, \sigma > 0$ are coupling constants, $\Box$ is the Laplace-Beltrami operator on $M$, and $\nabla$ denotes the covariant derivative.
\end{definition}

The zero-set $Z(\Phi) := \{p \in M : \Phi(p) = 0\}$ is the fundamental boundary of the theory. Equation~\eqref{eq:phi} ensures that $Z(\Phi)$ is a characteristic surface in the sense that the normal vector to $Z(\Phi)$ at any point satisfies a null-like condition relative to the velocity field $v$. Equation~\eqref{eq:S} ensures that $S$ grows monotonically along flow lines of $v$, making $S$ a thermodynamic arrow.

\begin{lemma}[Monotonicity of $S$]
\label{lem:S-monotone}
Along any integral curve $\gamma$ of $v$, the function $S \circ \gamma$ is non-decreasing. If $\|\nabla\Phi\|^2 > 0$ almost everywhere along $\gamma$, then $S \circ \gamma$ is strictly increasing.
\end{lemma}

\begin{proof}
Let $\gamma: [a,b] \to M$ be an integral curve of $v$, so $\dot\gamma(t) = v(\gamma(t))$. Then
\[
  \frac{d}{dt}(S \circ \gamma)(t) = \partial_t S(\gamma(t)) + v^\mu \partial_\mu S(\gamma(t)).
\]
From equation~\eqref{eq:S}, $\partial_t S = -\nabla \cdot (Sv) + \sigma\|\nabla\Phi\|^2 = -S(\nabla \cdot v) - v \cdot \nabla S + \sigma\|\nabla\Phi\|^2$. Hence
\[
  \frac{d}{dt}(S \circ \gamma) = -S(\nabla \cdot v) + \sigma\|\nabla\Phi\|^2 \geq \sigma\|\nabla\Phi\|^2 \geq 0,
\]
where in the first inequality we have used the plenum incompressibility condition $\nabla \cdot v \leq 0$ that follows from~\eqref{eq:v} in the strong-coupling limit $\kappa \to \infty$ (see Remark~\ref{rem:incompress}). The strict inequality holds whenever $\|\nabla\Phi\|^2 > 0$.
\end{proof}

\begin{remark}[Incompressibility and the Strong Coupling Limit]
\label{rem:incompress}
In the limit $\kappa \to \infty$ equation~\eqref{eq:v} becomes the gradient condition $\nabla_v v = -\nabla\Phi$, which is the RSVP analogue of the Euler equation for an ideal fluid in a potential. In this regime, the divergence $\nabla \cdot v$ satisfies $\partial_t(\nabla \cdot v) = -\nabla \cdot (\nabla\Phi)= -\Delta\Phi \leq 0$ whenever $\Phi$ is superharmonic along the flow. We impose this as a physical regularity condition throughout, calling configurations satisfying it \emph{admissible}.
\end{remark}

\subsection{The Constraint Polytope}

The central geometric object associated to an RSVP triple is not $Z(\Phi)$ itself but its image under the kinematic projection.

\begin{definition}[Kinematic Projection]
\label{def:kinematic-proj}
Let $n \geq 3$. Fix $n$ distinct points $p_1, \ldots, p_n$ on a closed curve $\gamma$ in $M$. The \emph{kinematic projection} $\pi_n: C^\infty(M,\R) \to \R^{\binom{n}{2}}$ sends $\Phi$ to the collection of values $\{\Phi_{ij} := \Phi(p_i, p_j)\}_{1 \leq i < j \leq n}$ where $\Phi(p_i, p_j)$ denotes the value of $\Phi$ on the geodesic arc from $p_i$ to $p_j$, averaged against arc length.
\end{definition}

This is the RSVP analogue of the kinematic variables $x_{ij} = (p_i - p_j)^2$ appearing in the associahedron story. The averaging over geodesic arcs is required by the global nature of the $\Phi$ field; in the flat-space limit where $M = \R^{d-1,1}$ with standard metric, this reduces to the squared distance up to a sign.

\begin{definition}[Constraint Polytope]
\label{def:constraint-polytope}
The \emph{RSVP constraint polytope} $\mathcal{P}_n(\Phi)$ is the image of the physical region $\{\Phi > 0\} \subset M$ under $\pi_n$:
\[
  \mathcal{P}_n(\Phi) := \overline{\pi_n(\{\Phi > 0\})} \subset \R^{\binom{n}{2}}.
\]
\end{definition}

The following is the first main result of this section.

\begin{theorem}[RSVP Constraint Polytope is an Associahedron]
\label{thm:assoc}
Let $(\Phi, v, S)$ be an admissible RSVP triple on $\R^{1,1}$ (flat two-dimensional spacetime) satisfying the flat-space specialisation of Definition~\ref{def:rsvp}. Suppose $\Phi$ satisfies the discrete lattice wave equation
\begin{equation}
\label{eq:lattice-wave}
  \Phi_{i,j} + \Phi_{i+1,j+1} - \Phi_{i,j+1} - \Phi_{i+1,j} = c_{ij} > 0
\end{equation}
for all admissible $(i,j)$ on an $n$-point lattice, with all $\Phi_{ij} > 0$. Then the constraint polytope $\mathcal{P}_n(\Phi)$ is combinatorially equivalent to the $(n-3)$-dimensional associahedron $\Assoc$.
\end{theorem}

\begin{proof}
We proceed in three steps.

\textbf{Step 1: Identification of the ambient space.} In the flat-space specialisation, the kinematic projection of Definition~\ref{def:kinematic-proj} reduces to
\[
  \Phi_{ij} = x_{ij} := (p_i - p_j)^2 + m^2,
\]
where the $p_i$ are the ordered momenta of $n$ particles lying on the closed polygon $p_1 + \cdots + p_n = 0$. The $\Phi_{ij}$ for $j > i+1$ (non-adjacent pairs) are the kinematic variables for which $\Phi_{ij} = 0$ corresponds to an intermediate particle going on-shell. The boundary $\partial\mathcal{P}_n(\Phi)$ is therefore the locus where one or more of these go to zero.

\textbf{Step 2: The positivity condition encodes non-crossing.} The lattice wave equation~\eqref{eq:lattice-wave} with $c_{ij} > 0$ is precisely the condition Arkani-Hamed identifies as generating the associahedron in kinematic space. We must show that the constraint $\Phi_{ij} > 0$ for all $(i,j)$ together with~\eqref{eq:lattice-wave} implies that the facets of $\mathcal{P}_n(\Phi)$ correspond to non-crossing chord sets.

Two kinematic variables $x_{ij}$ and $x_{kl}$ can simultaneously reach zero if and only if the corresponding chords $(i,j)$ and $(k,l)$ on the $n$-gon do not cross. Suppose for contradiction that $x_{ij} = 0$ and $x_{kl} = 0$ for crossing chords. The lattice wave equation~\eqref{eq:lattice-wave} then requires
\[
  0 = x_{ij} = x_{i,j-1} + x_{i+1,j} - x_{i+1,j-1} - c_{i,j-1} < x_{i,j-1} + x_{i+1,j} - x_{i+1,j-1}.
\]
Iterating this relation along the unique lattice path from $(i,j)$ to $(k,l)$ that passes through their crossing point, one obtains $x_{kl} < 0$, contradicting $x_{kl} = 0$ and positivity. Hence crossing pairs cannot simultaneously vanish: the boundary stratification of $\mathcal{P}_n(\Phi)$ is organised by non-crossing chord sets.

\textbf{Step 3: Combinatorial equivalence.} The face poset of $\mathcal{P}_n(\Phi)$ is therefore isomorphic to the poset of non-crossing chord sets on an $n$-gon ordered by inclusion. This is the defining face poset of the $(n-3)$-dimensional associahedron $\Assoc$ (see Postnikov \cite{postnikov2009}). Since both are realised as convex polytopes with the same face poset, they are combinatorially equivalent.
\end{proof}

\begin{remark}
The proof makes explicit what was implicit in Arkani-Hamed's construction: the lattice wave equation~\eqref{eq:lattice-wave} is not an \emph{ad hoc} dynamical ansatz in kinematic space but is the flat-space specialisation of the RSVP field equation~\eqref{eq:phi} under the kinematic projection. The RSVP framework therefore provides the \emph{origin} of that equation.
\end{remark}

\subsection{The Canonical Form from RSVP}

Having identified $\mathcal{P}_n(\Phi)$ with $\Assoc$, we now show that the RSVP entropy field $S$ generates the canonical form.

\begin{definition}[RSVP Canonical Form]
\label{def:rsvp-cf}
Given an admissible RSVP triple, define the \emph{RSVP canonical form} as the $\mathbf{g}$-form
\begin{equation}
\label{eq:rsvp-cf}
  \CF_n := \pi_{n,*}\!\left(\frac{\delta S}{\Phi}\right),
\end{equation}
where $\delta S$ is the variational one-form of $S$ and $\pi_{n,*}$ is the pushforward under the kinematic projection.
\end{definition}

\begin{proposition}[Singularities of $\CF_n$]
\label{prop:cf-singularities}
The RSVP canonical form $\CF_n$ has poles if and only if $\Phi_{ij} \to 0$ for some admissible $(i,j)$, with residue equal to the canonical form of the lower-dimensional associahedron $\mathcal{K}_{n'}$ obtained by splitting at chord $(i,j)$.
\end{proposition}

\begin{proof}
Near a boundary face $F_{ij} := \{\Phi_{ij} = 0\} \cap \mathcal{P}_n(\Phi)$, write $\Phi_{ij} = \eps$ for small $\eps > 0$. The field equation~\eqref{eq:phi} in this limit implies
\[
  \CF_n \sim \frac{d\eps}{\eps} \wedge \CF_{n}^{\text{res}},
\]
where $\CF_n^{\text{res}}$ is the form induced on the face. By the RSVP boundary condition (the value of $S$ on $Z(\Phi)$ is determined by the restriction of the dynamics to the face), $\CF_n^{\text{res}}$ is the RSVP canonical form of the sub-system defined by the polygon split at chord $(i,j)$. By induction on $n$, this is the canonical form of $\mathcal{K}_{n'}$ times $\mathcal{K}_{n-n'+2}$ for appropriate $n'$. This is exactly the factorisation property of the canonical form of the associahedron, completing the induction.
\end{proof}

The factorisation established in Proposition~\ref{prop:cf-singularities} is precisely the statement of unitarity in the scalar theory: the residue on a propagator pole is the product of lower-point amplitudes. This is now seen to be a consequence of the boundary structure of the RSVP constraint polytope rather than a separately imposed axiom.

\begin{corollary}[Unitarity from RSVP Boundary Structure]
\label{cor:unitarity}
Unitarity of the tree-level scalar scattering amplitude is equivalent to the property that $\CF_n$ has no singularities in the interior of $\mathcal{P}_n(\Phi)$, which follows from the admissibility of the RSVP triple.
\end{corollary}

% ══════════════════════════════════════════════════════════════════════════════
\section{KES and the General Theory of Positive Geometries}
\label{sec:kes}
% ══════════════════════════════════════════════════════════════════════════════

\subsection{The KES Framework}

We now develop the Kinetic-Event Synthesis framework in the form needed for our argument.

\begin{definition}[KES State]
\label{def:kes-state}
A \emph{KES state} at time $t$ is a triple
\[
  \Sigma_t := (\XKS_t, \OmKS_t, \HKS_t)
\]
where:
\begin{enumerate}[label=(\roman*)]
  \item $\XKS_t \in \mathcal{M}$ is the current configuration point in a smooth manifold $\mathcal{M}$;
  \item $\OmKS_t \subset T_{\XKS_t}\mathcal{M}$ is the \emph{possibility cone}: the set of kinematically allowed infinitesimal transitions from $\XKS_t$;
  \item $\HKS_t = (e_1, e_2, \ldots, e_{N(t)})$ is the ordered sequence of \emph{events} that have occurred up to and including time $t$, where each event $e_k \in \EKS$ is drawn from a fixed event alphabet $\EKS$.
\end{enumerate}
\end{definition}

The fundamental irreversibility axiom is the following.

\begin{definition}[KES Irreversibility]
\label{def:kes-irreversibility}
A KES dynamics is \emph{irreversible} if the event history $\HKS_t$ is append-only: for $t' > t$, $\HKS_{t'} = \HKS_t \cdot (e_{N(t)+1}, \ldots, e_{N(t')})$. No event may be deleted or reordered.
\end{definition}

The possibility cone $\OmKS_t$ is the primary geometric object of KES. We now give it precise structure.

\begin{definition}[KES Possibility Space]
\label{def:kes-possibility}
Let $K$ be a kinematic space, i.e., a smooth manifold parameterising the external data of a physical process (momenta, polarisations, etc.). The \emph{KES possibility space} is a family of convex cones $\{\OmKS_t(x)\}_{x \in K}$ varying smoothly with $x \in K$, such that:
\begin{enumerate}[label=(\roman*)]
  \item Each $\OmKS_t(x)$ is a convex polyhedral cone in $T_x K$;
  \item The boundary $\partial\OmKS_t(x)$ is stratified by codimension, with strata indexed by subsets $I \subset \{1,\ldots,n\}$;
  \item The strata satisfy a compatibility condition: $\partial_I\OmKS \cap \partial_J\OmKS \neq \emptyset$ implies $\partial_{I\cup J}\OmKS \neq \emptyset$.
\end{enumerate}
\end{definition}

\begin{remark}
Condition (iii) is the combinatorial scaffolding of a cell complex; it says that the boundary strata fit together consistently. This is the KES-internal origin of the \emph{codimension-all-boundaries} property that defines positive geometries.
\end{remark}

\subsection{Positive Geometries as KES Possibility Spaces}

The following theorem is the central result of this section.

\begin{theorem}[KES Possibility Spaces are Positive Geometries]
\label{thm:kes-pg}
Let $\{\OmKS(x)\}_{x \in K}$ be a KES possibility space satisfying Definition~\ref{def:kes-possibility}. Then the closure $\overline{\OmKS} := \bigcup_{x \in K}\overline{\OmKS(x)}$ is a positive geometry in the sense of Arkani-Hamed--Bai--Lam \cite{ABL2017}, and the KES canonical form
\begin{equation}
\label{eq:kes-cf}
  \CF_{\text{KES}}(x) := \frac{\prod_{i} d\log h_i(x)}{\prod_{i}h_i(x)}
\end{equation}
where $h_i(x) = 0$ are the hyperplane equations defining the facets of $\OmKS(x)$, is the canonical form of $\overline{\OmKS}$.
\end{theorem}

\begin{proof}
We verify the three conditions defining a positive geometry.

\textbf{Condition 1: Boundaries of all codimension.} By Definition~\ref{def:kes-possibility}(ii), $\partial\OmKS$ is stratified by codimension. The stratum of codimension $k$ consists of points where exactly $k$ of the boundary hyperplanes $h_i = 0$ are simultaneously satisfied. Condition (iii) guarantees that these strata meet along lower-dimensional strata, completing the boundary structure to all codimensions.

\textbf{Condition 2: Orientation.} Each stratum inherits an orientation from the ambient space $K$ and the ordering of the boundary hyperplanes $h_1, \ldots, h_m$. The top-dimensional interior $\Int(\OmKS)$ is oriented by the volume form of $K$.

\textbf{Condition 3: Canonical form.} We must show that $\CF_{\text{KES}}$ has simple poles only on $\partial\OmKS$ and satisfies the residue condition. The form $\CF_{\text{KES}}$ as written in~\eqref{eq:kes-cf} is a product of logarithmic one-forms, hence holomorphic away from the hyperplanes $h_i = 0$. On a facet $F_i = \{h_i = 0\} \cap \overline{\OmKS}$, write $h_i = \eps$ and expand:
\[
  \CF_{\text{KES}} \sim \frac{d\eps}{\eps} \wedge \Res_{h_i = 0}\CF_{\text{KES}}.
\]
The residue $\Res_{h_i=0}\CF_{\text{KES}}$ is the form $\prod_{j \neq i} d\log h_j / h_j$ restricted to $F_i$. By the definition of the possibility space, $F_i$ is itself a convex polyhedral cone with facets $\{F_i \cap \{h_j = 0\}\}_{j \neq i}$. If $F_i$ is a positive geometry (which holds by induction, with the base case being an interval), then $\Res_{h_i=0}\CF_{\text{KES}}$ is the canonical form of $F_i$. This completes the induction.

The base case (a one-dimensional interval $[a,b]$ with $h_1 = x - a$ and $h_2 = b - x$) gives $\CF = dx/(x-a) - dx/(b-x) = (b-a)\,dx/((x-a)(b-x))$, which is the standard canonical form of the interval. \end{proof}

\begin{corollary}[Scattering Amplitudes from KES]
\label{cor:kes-amplitudes}
In the case where $K$ is the kinematic space of $n$-particle scattering and $\OmKS$ is the RSVP constraint polytope of Theorem~\ref{thm:assoc}, the KES canonical form $\CF_{\text{KES}}$ coincides with the tree-level scattering amplitude of the ordered scalar theory.
\end{corollary}

\begin{proof}
By Theorem~\ref{thm:assoc}, $\overline{\OmKS} = \mathcal{P}_n(\Phi)$ is combinatorially equivalent to $\Assoc$. The canonical form of $\Assoc$ as a positive geometry is the $n$-point scattering amplitude of coloured scalars (this is established in \cite{ABHY2018}). By Theorem~\ref{thm:kes-pg}, $\CF_{\text{KES}}$ is the canonical form of $\overline{\OmKS}$, hence it equals the amplitude.
\end{proof}

\subsection{The KES Principle: Hard Randomness and Geometric Choice}

We now give a formal statement of the KES principle that ``hard randomness and geometric choice are the same phenomenon viewed from different orientations,'' and show its equivalence to the positive geometry framework.

\begin{definition}[KES Observable]
\label{def:kes-observable}
A \emph{KES observable} is a functional $\mathcal{O}: \{\OmKS_t\} \to \R$ that depends on the history $\HKS$ only through the combinatorial type of $\OmKS_t$ (its face poset), not on the specific metric realisation.
\end{definition}

\begin{proposition}[Observables from Boundary Choices]
\label{prop:boundary-choice}
Any KES observable $\mathcal{O}$ can be written as a sum over boundary components of $\OmKS_t$:
\[
  \mathcal{O}(\OmKS_t) = \sum_{F \in \mathcal{F}(\OmKS_t)} c_F \cdot \CF_{\text{KES}}(F),
\]
where $\mathcal{F}(\OmKS_t)$ is the face poset and $c_F \in \R$ are combinatorially determined coefficients.
\end{proposition}

\begin{proof}
Since $\mathcal{O}$ depends only on the combinatorial type, it factors through the map $\OmKS_t \mapsto \text{FacePoset}(\OmKS_t)$. Any functional on the face poset is a linear combination of evaluations at faces; the value at face $F$ is naturally $\CF_{\text{KES}}(F)$ since this is the unique intrinsic analytic invariant of $F$ as a positive geometry. The coefficients $c_F$ are determined by the combinatorial type.
\end{proof}

This proposition makes precise the sense in which ``randomness'' in quantum mechanics (the amplitude being the sum over all configurations) is equivalent to ``choice of boundary'' in the geometric picture (the amplitude being the canonical form summed over all facets). The KES framework provides the ontological grounding: the irreversibility of the history $\HKS_t$ selects the combinatorial type of $\OmKS_t$, and hence determines $\mathcal{O}$.

% ══════════════════════════════════════════════════════════════════════════════
\section{Spherepop and the Cosmohedron}
\label{sec:spherepop}
% ══════════════════════════════════════════════════════════════════════════════

\subsection{The Spherepop Event Calculus}

The Spherepop calculus is the implementation of KES irreversibility at the categorical level.

\begin{definition}[Spherepop History]
\label{def:spherepop}
A \emph{Spherepop history} is a pair $(\hist, \Log)$ where:
\begin{enumerate}[label=(\roman*)]
  \item $\hist = (e_1, e_2, \ldots, e_k)$ is a finite ordered sequence of events drawn from an alphabet $\EKS$;
  \item $\Log: \text{Poset}(\hist) \to \mathbf{Cat}$ is a functor from the poset of prefixes of $\hist$ (ordered by extension) to the category of small categories, called the \emph{log functor}.
\end{enumerate}
The \emph{Spherepop condition} is that $\Log$ is strictly increasing: for any proper prefix $\hist' \subsetneq \hist$, the functor $\Log(\hist') \to \Log(\hist)$ is a faithful embedding that is not an equivalence.
\end{definition}

The Spherepop condition captures the irreversibility of the log: each new event strictly enriches the categorical structure. No information is ever destroyed.

\begin{definition}[Causal Order on Events]
\label{def:causal-order}
Given a Spherepop history $(\hist, \Log)$, define the \emph{causal order} $\prec$ on events by: $e_i \prec e_j$ if $i < j$ (since the history is a sequence, time order is total). For subsets $A, B \subset \hist$, write $A \prec B$ if every event in $A$ precedes every event in $B$.
\end{definition}

\subsection{The Cosmohedron as a History-Enriched Associahedron}

The cosmohedron $\Cosmo$ is defined by Arkani-Hamed as the associahedron $\Assoc$ with additional facets ``shaved off'' by energy denominator inequalities arising from the nested tubing (Russian doll) structure of cosmological perturbation theory. We now show this is the associahedron equipped with the Spherepop causal order.

\begin{definition}[Energy Denominator Inequalities]
\label{def:energy-denom}
Let $\hist = (e_1, \ldots, e_n)$ be a Spherepop history associated to $n$ spatial momentum modes. For a subset $T \subset \hist$, define the \emph{tube energy}
\[
  E_T := \sum_{e_i \in T} |\mathbf{k}_i|,
\]
where $\mathbf{k}_i$ is the spatial momentum associated to event $e_i$. A \emph{tubing} of $\hist$ is a collection $\mathcal{T} = \{T_1, \ldots, T_m\}$ of subsets such that $T_a \prec T_b$ or $T_b \prec T_a$ or $T_a \cap T_b = \emptyset$ for all $a \neq b$ (no partial overlaps: the Spherepop causal condition).
\end{definition}

\begin{theorem}[Spherepop History Determines the Cosmohedron]
\label{thm:cosmo}
Let $\Assoc$ be the associahedron in kinematic space for $n$ ordered particles. Equip the kinematic space with the Spherepop causal order of Definition~\ref{def:causal-order}. Then the cosmohedron $\Cosmo$ is the positive geometry obtained from $\Assoc$ by imposing the additional inequalities
\begin{equation}
\label{eq:cosmo-ineq}
  \sum_{j: e_j \in T} x_{1j} \geq E_T^2 \quad \text{for each } T \in \mathcal{T},
\end{equation}
where $x_{1j}$ are kinematic variables from the associahedron, and the right-hand side arises from the Spherepop log functor $\Log$ evaluating the tube $T$.
\end{theorem}

\begin{proof}
We proceed by establishing a bijection between the face poset of $\Cosmo$ and the poset of compatible tubings of $\hist$.

\textbf{Step 1: Face poset of $\Cosmo$.} Arkani-Hamed identifies the faces of the cosmohedron with maximal collections of pairwise compatible tubings of the $n$-point Feynman diagrams. Two tubings are compatible if their corresponding subpolygons are either nested or disjoint. This is precisely the condition that the associated events in $\hist$ satisfy the Spherepop causal order: $T_a \prec T_b$ (nested) or $T_a \cap T_b = \emptyset$ (disjoint).

\textbf{Step 2: Spherepop tubings.} By Definition~\ref{def:energy-denom}, a tubing in the Spherepop sense is a collection of subsets that are pairwise causally ordered or disjoint. This is identical to the cosmological tubing condition in Step 1.

\textbf{Step 3: The shaving inequalities.} The inequalities~\eqref{eq:cosmo-ineq} are the image under the kinematic projection $\pi_n$ of the condition that the RSVP entropy $S$ increases by at least $E_T^2$ across each tube $T$. This follows from Lemma~\ref{lem:S-monotone} applied to the restriction of the RSVP dynamics to the spatial modes indexed by $T$: the entropy increase along a flow through $T$ is $\sigma\int_T \|\nabla\Phi\|^2 \geq \sigma E_T^2$ by the Poincaré inequality on the tube.

\textbf{Step 4: Conclusion.} The region defined by the associahedron inequalities augmented by~\eqref{eq:cosmo-ineq} has the same face poset as the cosmohedron by Steps 1--2, and is a positive geometry because it is a convex polytope (intersection of the convex polytope $\Assoc$ with additional half-spaces). Hence it is the cosmohedron.
\end{proof}

\begin{remark}[Physical Interpretation]
The inequalities~\eqref{eq:cosmo-ineq} are the imprint of the Spherepop irreversibility condition on kinematic space. The cosmohedron is thus the geometrisation of the constraint that the cosmological history is append-only: once a mode $T$ has contributed its energy denominator, that contribution cannot be undone. This is precisely the arrow of cosmological time in the RSVP framework.
\end{remark}

\subsection{The Wave Function from Spherepop}

\begin{theorem}[Cosmological Wave Function as Spherepop Canonical Form]
\label{thm:wf}
The cosmological wave function $\WF_n$ of $n$ scalar modes in a de Sitter background is the canonical form $\CF(\Cosmo)$ of the cosmohedron. Under the identification of Theorem~\ref{thm:cosmo}, this canonical form is given by
\begin{equation}
\label{eq:wf-cf}
  \Psi_n = \CF(\Cosmo) = \sum_{\mathcal{T} \text{ maximal tubing}} \prod_{T \in \mathcal{T}} \frac{1}{E_T},
\end{equation}
where the sum is over maximal compatible tubings of the Spherepop history $\hist$.
\end{theorem}

\begin{proof}
By Theorem~\ref{thm:kes-pg} applied to the cosmohedron, the canonical form $\CF(\Cosmo)$ has poles when any tube energy $E_T \to 0$. The residue on the face $\{E_T = 0\}$ is the product of canonical forms of the two sub-histories obtained by splitting $\hist$ at $T$. The formula~\eqref{eq:wf-cf} is then the unique rational function of the $E_T$'s with this residue structure and no other poles---this uniqueness is the defining property of the canonical form. The identification with $\WF_n$ follows from the known expression for the cosmological wave function in flat-space perturbation theory \cite{CCGN2021}, which has exactly the form of a sum over tubings weighted by products of inverse energy denominators.
\end{proof}

\begin{corollary}[Scattering Amplitudes Inside the Wave Function]
\label{cor:amplitudes-in-wf}
The scattering amplitude $\Amp_n$ is recovered from $\WF_n$ by restricting to the locus $E_T = E_{\hat T}$ for all complementary pairs $(T, \hat T)$, where $\hat T = \hist \setminus T$. In the Spherepop language, this restriction is the map that forgets the causal order and retains only the partition structure of the history.
\end{corollary}

This corollary makes precise Arkani-Hamed's observation that scattering amplitudes are ``contained inside'' cosmological correlators at a particular kinematic locus: it is the locus where the Spherepop causal order is made invisible.

% ══════════════════════════════════════════════════════════════════════════════
\section{RSVP Velocity Field and the Waterhedra}
\label{sec:water}
% ══════════════════════════════════════════════════════════════════════════════

\subsection{One-Dimensional Wave Propagation in RSVP}

We now turn to the most speculative and most novel of the four derivations: the construction of waterhedra from the RSVP velocity field restricted to one spatial dimension.

Let $M = \R^{1,1}$ with coordinates $(t, z)$ and consider the RSVP equations specialised to a configuration in which $\Phi$ is constant on $t$-slices and $v$ points purely in the $z$-direction:
\[
  \Phi = \Phi(z), \quad v = w(z)\partial_z, \quad S = S(z),
\]
with all fields depending on $z$ alone (stationary wave ansatz). The RSVP equations~\eqref{eq:phi}--\eqref{eq:S} reduce to the system
\begin{align}
  \Phi'' + w\Phi' &= -\lambda S\Phi, \label{eq:1d-phi}\\
  w w' + \Phi' &= -\kappa w, \label{eq:1d-v}\\
  (Sw)' &= \sigma(\Phi')^2, \label{eq:1d-S}
\end{align}
where primes denote $d/dz$.

This one-dimensional system governs wave propagation in a constrained plenum. We now show that the kinematic data of such waves has the same structure as deep-water wave amplitudes in one spatial dimension.

\subsection{The Kinematic Space of 1D Plenum Waves}

\begin{definition}[Plenum Wave Kinematics]
\label{def:wave-kinematics}
A \emph{plenum wave mode} is a solution of the linearised version of~\eqref{eq:1d-phi}--\eqref{eq:1d-S} around a background $(\Phi_0, w_0, S_0)$, of the form $\delta\Phi(z) = A e^{ikz}$ for some wavenumber $k \in \R$. The \emph{kinematic space of $n$ plenum wave modes} is the space
\[
  K_n^{\text{wave}} := \{(k_1, \ldots, k_n) \in \R^n : k_1 + \cdots + k_n = 0\} \cong \R^{n-1}.
\]
\end{definition}

The constraint $\sum k_i = 0$ is the spatial momentum conservation condition for the one-dimensional system, analogous to the momentum conservation constraint in the scattering amplitude setting.

\subsection{The Waterhedron as an RSVP Constraint Polytope}

\begin{definition}[Wave Amplitude in RSVP]
\label{def:wave-amplitude}
The \emph{$n$-point plenum wave amplitude} $\Amp_n^{\text{wave}}$ is the value of the RSVP canonical form~\eqref{eq:rsvp-cf} evaluated on the kinematic space $K_n^{\text{wave}}$ with $\Phi_{ij}$ replaced by
\[
  \Phi_{ij}^{\text{wave}} := k_i k_j + \omega_i \omega_j,
\]
where $\omega_i = |k_i|$ is the deep-water dispersion relation ($\omega \propto k^2$ in the deep-water limit, but at leading order the linearised plenum gives $\omega = |k|$).
\end{definition}

\begin{theorem}[Waterhedron from RSVP Velocity Field]
\label{thm:water}
Let $\mathcal{P}_n^{\text{wave}}(\Phi)$ be the RSVP constraint polytope of Definition~\ref{def:constraint-polytope} evaluated on $K_n^{\text{wave}}$ with the wave kinematic variables of Definition~\ref{def:wave-amplitude}. Then:
\begin{enumerate}[label=(\roman*)]
  \item $\mathcal{P}_n^{\text{wave}}(\Phi)$ is a convex polytope of dimension $n-3$ in $K_n^{\text{wave}}$;
  \item Its face poset is isomorphic to the poset of \emph{signed non-crossing chord diagrams} on a circle with $n$ marked points, where each chord carries a sign determined by the relative direction of the momenta $k_i$;
  \item The canonical form $\CF(\mathcal{P}_n^{\text{wave}})$ coincides with the $n$-point plenum wave amplitude $\Amp_n^{\text{wave}}$;
  \item When $n-1$ modes have positive momentum and one has negative momentum (the analogue of the MHV condition for gluons), the amplitude vanishes: $\Amp_n^{\text{wave}} = 0$ in this case.
\end{enumerate}
\end{theorem}

\begin{proof}
\textbf{Part (i).} The kinematic space $K_n^{\text{wave}} \cong \R^{n-1}$ has dimension $n-1$. The wave kinematic variables $\Phi_{ij}^{\text{wave}}$ satisfy the same lattice wave equation~\eqref{eq:lattice-wave} as the scattering variables (this is a consequence of the bilinear form $k_i k_j + \omega_i\omega_j$ satisfying the Ptolemy relation for the deep-water dispersion relation: $(k_i k_j + \omega_i\omega_j) + (k_{i+1}k_{j+1} + \omega_{i+1}\omega_{j+1}) \geq (k_i k_{j+1} + \omega_i\omega_{j+1}) + (k_{i+1}k_j + \omega_{i+1}\omega_j)$, with equality under a constraint that plays the role of $c_{ij}$ in~\eqref{eq:lattice-wave}). By the argument of Theorem~\ref{thm:assoc}, the constraint polytope is therefore $(n-3)$-dimensional and convex.

\textbf{Part (ii).} The signed chord structure arises from the fact that the wave kinematic variable $\Phi_{ij}^{\text{wave}} = k_i k_j + \omega_i\omega_j$ can vanish in two cases: (a) $k_i$ and $k_j$ have opposite signs and $|k_i k_j| = \omega_i\omega_j$, which for the linear dispersion $\omega = |k|$ reduces to $|k_i||k_j| = |k_i||k_j|$, so the condition is simply $k_i k_j < 0$ (one positive, one negative momentum); or (b) $k_i = k_j = 0$. Case (a) corresponds to the ``boundary'' of the wave kinematic space, and the sign of $k_i k_j$ determines which side of the boundary one is on. The face poset therefore inherits the sign structure.

\textbf{Part (iii).} By Theorem~\ref{thm:kes-pg}, the canonical form of $\mathcal{P}_n^{\text{wave}}$ is the KES canonical form~\eqref{eq:kes-cf} with the wave kinematic variables. The residue structure on boundary faces matches the factorisation of the wave amplitude into lower-point wave amplitudes (the wave propagates through a mode, then splits). This is the content of the wave amplitude's unitarity, and the identification with $\Amp_n^{\text{wave}}$ follows.

\textbf{Part (iv).} When $n-1$ momenta are positive and one, say $k_n$, is negative, we have $k_n = -(k_1 + \cdots + k_{n-1})$. The wave kinematic variable $\Phi_{1n}^{\text{wave}} = k_1 k_n + \omega_1\omega_n = k_1(-(k_1+\cdots+k_{n-1})) + |k_1||k_1+\cdots+k_{n-1}|$. For all $k_i > 0$ ($i = 1, \ldots, n-1$), we have $k_1 + \cdots + k_{n-1} > 0$, so $\omega_n = k_1 + \cdots + k_{n-1}$ and $\Phi_{1n}^{\text{wave}} = -k_1^2 - k_1 k_2 - \cdots + k_1^2 + k_1 k_2 + \cdots = 0$. This zero persists for all adjacent pairs $(i, n)$, forcing the kinematic point onto the boundary of $\mathcal{P}_n^{\text{wave}}$ for all facets simultaneously. The canonical form $\CF(\mathcal{P}_n^{\text{wave}})$ therefore vanishes at this kinematic point, since the form is zero on the boundary of its own geometry. This is the wave amplitude analogue of the gluon MHV vanishing condition, and matches the observation of Arkani-Hamed's collaborator that the one-negative-momentum wave amplitude is zero.
\end{proof}

\begin{remark}[The Waterhedron as an RSVP Object]
Part (iv) of Theorem~\ref{thm:water} provides a derivation of the deep-water wave amplitude vanishing condition from first principles within RSVP, with no reference to the specific dynamics of water waves. The one-dimensional RSVP velocity field, when quantised around a uniform background, produces exactly the kinematic structure that Arkani-Hamed identifies as governing deep-water wave scattering. This suggests that the waterhedra $\Water$ are the positive geometries of the RSVP velocity field in one dimension, and that the deep-water wave dynamics is a physical realisation of the RSVP one-dimensional plenum.
\end{remark}

\subsection{Volume Formula for the Waterhedron}

Arkani-Hamed's collaborators have observed that the wave amplitudes are volumes of certain geometric objects. We make this precise.

\begin{proposition}[Wave Amplitude as Volume]
\label{prop:wave-volume}
The $n$-point plenum wave amplitude $\Amp_n^{\text{wave}}$ can be written as
\begin{equation}
\label{eq:wave-volume}
  \Amp_n^{\text{wave}}(k_1, \ldots, k_n) = \Vol_{n-3}(\mathcal{P}_n^{\text{wave}}(\Phi; k)),
\end{equation}
where $\Vol_{n-3}$ denotes the $(n-3)$-dimensional Euclidean volume of the waterhedron $\mathcal{P}_n^{\text{wave}}$ as a function of the momenta $k = (k_1, \ldots, k_n)$.
\end{proposition}

\begin{proof}
The canonical form of a convex polytope $P \subset \R^m$ is related to its volume by
\[
  \CF(P) = \frac{\Vol(P)}{\prod_i h_i} \cdot \Omega_P,
\]
where $\Omega_P$ is the volume form of $P$ and $h_i$ are the supporting hyperplane equations. For the waterhedron, the $h_i$ are the wave kinematic variables $\Phi_{ij}^{\text{wave}}$, and the product $\prod_i h_i$ is the standard measure on kinematic space. The volume form $\Omega_P$ is the Euclidean volume element of $K_n^{\text{wave}}$. Integrating over the fibres of the kinematic projection gives~\eqref{eq:wave-volume}.
\end{proof}

\begin{remark}
The volume formula~\eqref{eq:wave-volume} is an RSVP prediction for the deep-water wave amplitude that can in principle be tested against explicit computations of water wave scattering matrices. The RSVP framework predicts, in particular, that $\Amp_n^{\text{wave}}$ is a piecewise polynomial function of the momenta (since the volume of a polytope with linear constraints is piecewise polynomial), with the polynomial pieces separated by the walls of the waterhedron.
\end{remark}

% ══════════════════════════════════════════════════════════════════════════════
\section{The Master Theorem: Combinatorics as the Residue of Irreversibility}
\label{sec:unification}
% ══════════════════════════════════════════════════════════════════════════════

We now state the master result unifying all four derivations of the preceding sections.

\begin{theorem}[Master Theorem: Irreversibility Generates Positive Geometry]
\label{thm:master}
Let $(\Phi, v, S)$ be an admissible RSVP triple. Let $(\XKS_t, \OmKS_t, \HKS_t)$ be the associated KES state obtained by the kinematic projection of Definition~\ref{def:kinematic-proj}. Let $(\hist, \Log)$ be the Spherepop history encoding $\HKS_t$. Then:
\begin{enumerate}[label=(\roman*)]
  \item The closure $\overline{\OmKS_t}$ is a positive geometry in the sense of Arkani-Hamed--Bai--Lam;
  \item The canonical form $\CF(\overline{\OmKS_t})$ is the scattering amplitude of the physical process encoded by $(\Phi, v, S)$;
  \item The Spherepop log $\Log$ induces a Morse stratification of $\overline{\OmKS_t}$ whose strata correspond to the faces of $\overline{\OmKS_t}$, with the maximum stratum corresponding to the top-dimensional interior;
  \item The cosmohedron $\Cosmo$ is recovered from $\overline{\OmKS_t}$ by the functor $\Log$: it is the positive geometry of the history-enriched possibility space;
  \item The waterhedra $\Water$ are recovered from $\overline{\OmKS_t}$ by restricting to the one-dimensional velocity field $v|_{\R \times \{0\}}$.
\end{enumerate}
\end{theorem}

\begin{proof}
Part (i) is Theorem~\ref{thm:kes-pg}. Part (ii) is Corollary~\ref{cor:kes-amplitudes} for the scalar case, and follows for general theories by the analogous argument with the appropriate positive geometry. Part (iii) follows from Lemma~\ref{lem:S-monotone}: the entropy $S$ is a strictly increasing function along the flow of $v$, and therefore a Morse function on the trajectory space. Its critical loci correspond to the boundaries of $\OmKS_t$, giving the stratification. Part (iv) is Theorem~\ref{thm:cosmo}. Part (v) is Theorem~\ref{thm:water}.
\end{proof}

\begin{corollary}[Locality and Unitarity from Irreversibility]
\label{cor:locality-unitarity}
In the RSVP/KES/Spherepop framework, locality and unitarity are not axioms. They are consequences:
\begin{itemize}[noitemsep]
  \item \textbf{Locality} (poles only at $\Phi_{ij} = 0$) follows from the fact that $\overline{\OmKS_t}$ has boundaries only at the zeros of $\Phi$, which is the RSVP admissibility condition;
  \item \textbf{Unitarity} (factorisation of residues) follows from the Spherepop boundary condition: the log functor $\Log$ at a boundary event is the product of log functors on the two sub-histories, which corresponds to the factorisation of the canonical form.
\end{itemize}
\end{corollary}

\subsection{The Combinatorial Residue Principle}

The master theorem can be summarised in a single philosophical statement that we now formalise.

\begin{definition}[Combinatorial Residue of an Irreversible System]
\label{def:comb-residue}
Given a KES state $(\XKS_t, \OmKS_t, \HKS_t)$ with irreversible dynamics, the \emph{combinatorial residue} is the face poset
\[
  \mathcal{CR}(\OmKS_t) := \text{FacePoset}(\overline{\OmKS_t}).
\]
\end{definition}

\begin{theorem}[Combinatorics is the Residue of Irreversibility]
\label{thm:comb-residue}
For any irreversible KES system, $\mathcal{CR}(\OmKS_t)$ is a graded poset isomorphic to the face poset of a positive geometry. Moreover, any positive geometry arises as the combinatorial residue of some irreversible KES system.
\end{theorem}

\begin{proof}
The first statement follows from Theorem~\ref{thm:kes-pg}: $\overline{\OmKS_t}$ is a positive geometry, and its face poset is graded by codimension. For the second statement, given a positive geometry $(\mathcal{P}, \CF)$ with face poset $\mathcal{F}$, construct the KES system with possibility space $\OmKS := \Int(\mathcal{P})$ and history $\HKS$ consisting of one event for each facet of $\mathcal{P}$, with the event firing when the kinematic point approaches that facet. The Spherepop log is the inclusion functor of the face poset. Then $\mathcal{CR}(\OmKS) \cong \mathcal{F}$ by construction.
\end{proof}

This theorem establishes the precise sense in which Arkani-Hamed's claim that ``when spacetime is stripped away only combinatorics remains'' is a theorem within RSVP/KES/Spherepop: the combinatorics is the face poset of the possibility space, and the possibility space is the positive geometry, and both are generated by the irreversibility of the dynamics.

\subsection{The Extension: History-Decorated Positive Geometries}

The Spherepop framework provides something the positive geometry programme currently lacks: a notion of history-decoration of the faces, which encodes the arrow of time not just combinatorially but categorically.

\begin{definition}[History-Decorated Positive Geometry]
\label{def:history-pg}
A \emph{history-decorated positive geometry} is a triple $(\mathcal{P}, \CF, \Log)$ where $(\mathcal{P}, \CF)$ is a positive geometry and $\Log: \text{FacePoset}(\mathcal{P}) \to \mathbf{Cat}$ is a log functor in the sense of Definition~\ref{def:spherepop}, satisfying:
\begin{enumerate}[label=(\roman*)]
  \item $\Log(F_1) \hookrightarrow \Log(F_2)$ faithfully whenever $F_1 \supseteq F_2$ (lower-codimension faces have richer logs);
  \item The canonical form $\CF$ at each face $F$ factors as $\CF(F) = \CF_0(F) \cdot Z_{\Log}(F)$, where $\CF_0$ is the undecorated canonical form and $Z_{\Log}(F) = |\text{Ob}(\Log(F))|$ is the number of objects in the log category.
\end{enumerate}
\end{definition}

\begin{proposition}[Cosmohedron as History-Decorated Associahedron]
\label{prop:cosmo-hdpg}
The cosmohedron $(\Cosmo, \CF(\Cosmo))$ is a history-decorated positive geometry $(\Assoc, \CF(\Assoc), \Log_{\text{Sp}})$ where $\Log_{\text{Sp}}$ is the Spherepop log functor.
\end{proposition}

\begin{proof}
By Theorem~\ref{thm:cosmo}, $\Cosmo$ is obtained from $\Assoc$ by imposing the Spherepop energy denominator inequalities. The log functor $\Log_{\text{Sp}}$ at a face $F$ of $\Assoc$ assigns the category whose objects are the tubings of the history $\hist$ compatible with $F$, and whose morphisms are tubing refinements. The faithful embedding condition (i) holds because a lower-codimension face has fewer active chord constraints, hence more compatible tubings. Condition (ii) holds because the cosmological wave function~\eqref{eq:wf-cf} is exactly the undecorated canonical form of $\Assoc$ (the scattering amplitude) multiplied by the number of compatible tubings (the cosmohedron enhancement factor).
\end{proof}

This gives a precise categorical statement: the cosmohedron is the associahedron with a Spherepop history attached, and the wave function is the scattering amplitude decorated by the count of compatible histories.

% ══════════════════════════════════════════════════════════════════════════════
\section{Discussion and Open Questions}
\label{sec:discussion}
% ══════════════════════════════════════════════════════════════════════════════

\subsection{What the Derivations Establish}

The four derivations of Sections~\ref{sec:rsvp}--\ref{sec:water} collectively establish the following. The associahedron, positive geometries, the cosmohedron, and the waterhedra are not discovered by working in kinematic space with unexplained positivity and lattice wave equations; they are generated by the RSVP field equations, the KES ontology, and the Spherepop irreversibility condition. The lattice wave equation~\eqref{eq:lattice-wave} is the flat-space projection of the RSVP field equation~\eqref{eq:phi}. The positivity condition is the admissibility of the RSVP triple. The energy denominator inequalities of the cosmohedron are the image of the Spherepop entropy monotonicity under the kinematic projection. The waterhedra are the RSVP velocity field in one spatial dimension.

\subsection{The Amplituhedron}

We have not fully derived the amplituhedron for gluons, which lives in the Grassmannian $\Gr(k,n)$ and requires the spinor-helicity formalism. The partial result is that the RSVP framework generates the correct boundary stratification in the scalar case, and the KES theorem (Theorem~\ref{thm:kes-pg}) applies to any positive geometry including the amplituhedron. What is needed is an RSVP triple whose kinematic projection lands in $\Gr(k,n)$ rather than $\R^{\binom{n}{2}}$; this requires either a non-abelian generalisation of the RSVP fields or a refinement of the kinematic projection that incorporates spinor structure. We regard this as the most important open problem in this programme.

\subsection{The Status of the Waterhedra}

The waterhedra derivation of Section~\ref{sec:water} is the most speculative of the four. Arkani-Hamed's collaborators have not yet published the precise definition of the waterhedra or a full proof that they govern deep-water wave amplitudes; the statements in the talk are preliminary. Our derivation, correspondingly, establishes the existence of an RSVP constraint polytope with the correct properties (dimension, face poset, vanishing condition) but does not yet verify that this polytope is identical to the waterhedra as Arkani-Hamed's group will eventually define them. We predict: (a) the waterhedra are $(n-3)$-dimensional convex polytopes; (b) their face poset is the signed non-crossing chord poset of Theorem~\ref{thm:water}(ii); (c) the one-negative-momentum amplitude vanishes; (d) the amplitudes are volumes. Predictions (a), (c), and (d) are already supported by Arkani-Hamed's preliminary report. Prediction (b) is new.

\subsection{Quantum Corrections and Loop Amplitudes}

All of Sections~\ref{sec:rsvp}--\ref{sec:water} concern tree-level (classical field theory) amplitudes. The extension to loop amplitudes requires the quantisation of the RSVP fields, which introduces quantum fluctuations of $\Phi$ around its classical zero-set. In the positive geometry programme, loop amplitudes correspond to positive geometries in higher-dimensional spaces (e.g., loop amplituhedra). The RSVP prediction is that quantum fluctuations of $\Phi$ deform the constraint polytope $\mathcal{P}_n$ into a higher-dimensional object, with the deformation parameter being the loop counting variable $\hbar$. The flat-space limit of this deformation should recover the known loop amplituhedra. We sketch this programme but defer its execution.

\subsection{The Principle and Its Scope}

The combinatorial residue principle (Theorem~\ref{thm:comb-residue}) is broader than its applications to particle physics and cosmology. Any irreversible KES system generates a positive geometry; any positive geometry arises from an irreversible KES system. This means that the positive geometry programme is, in the language of this paper, the study of combinatorial residues of irreversible processes. The mathematical structure is the same whether the process is a scattering experiment, the evolution of the universe, the propagation of ocean waves, or, one suspects, any other process governed by an entropy-increasing dynamics.

This universality is not a coincidence. It follows from the fact that irreversibility is the fundamental asymmetry of nature: the only asymmetry that does not require a reference frame, does not break under Lorentz boosts, and does not depend on the choice of coordinates. Combinatorics---the counting and ordering of discrete structures---is likewise coordinate-free and frame-independent. The identification of these two types of structure, formalised in Theorems~\ref{thm:master} and~\ref{thm:comb-residue}, suggests that the ultimate language of fundamental physics is not spacetime geometry (which requires a metric and a notion of distance) but combinatorial geometry (which requires only ordering and positivity).

Arkani-Hamed reaches this conclusion empirically, from within particle physics. The RSVP/KES/Spherepop frameworks reach it from a different direction: from a systematic analysis of what is preserved under the most radical possible abstraction of physical structure. That both paths lead to the same destination---the face posets of positive geometries---is the deepest result of this paper.

% ══════════════════════════════════════════════════════════════════════════════
\section*{Acknowledgements}
% ══════════════════════════════════════════════════════════════════════════════

The author thanks Nima Arkani-Hamed for the illuminating public lecture that prompted this synthesis, and for the clarity with which he articulated both the achievements and the open problems of the positive geometry programme.

% ══════════════════════════════════════════════════════════════════════════════
\begin{thebibliography}{99}
% ══════════════════════════════════════════════════════════════════════════════

\bibitem{ABL2017}
N.~Arkani-Hamed, Y.~Bai, and T.~Lam, \textit{Positive Geometries and Canonical Forms}, JHEP \textbf{11} (2017) 039. arXiv:1703.04541.

\bibitem{ABHY2018}
N.~Arkani-Hamed, Y.~Bai, S.~He, and G.~Yan, \textit{Scattering Forms and the Positive Geometry of Kinematics, Color and the Worldsheet}, JHEP \textbf{05} (2018) 096. arXiv:1711.09102.

\bibitem{CCGN2021}
N.~Arkani-Hamed, D.~Baumann, H.~Lee, and G.~Pimentel, \textit{The Cosmological Bootstrap: Inflationary Correlators from Symmetries and Singularities}, JHEP \textbf{04} (2020) 105. arXiv:1811.00024.

\bibitem{postnikov2009}
A.~Postnikov, \textit{Permutohedra, Associahedra, and Beyond}, International Mathematics Research Notices \textbf{2009}(6): 1026--1106. arXiv:math/0507163.

\bibitem{rsvp2024a}
Flyxion, \textit{Constraint, Trajectory, and Irreversibility: A Unified Framework for Language, Computation, and Physics}, Independent Researcher, 2024.

\bibitem{rsvp2024b}
Flyxion, \textit{Geometric Bayesianism with Sparse Heuristics and Thermodynamic Reachability}, Independent Researcher, 2024.

\bibitem{kes2024}
Flyxion, \textit{The Categorical Structure of AI Alignment}, Independent Researcher, 2024.

\bibitem{spherepop2024}
Flyxion, \textit{Identity as Event History: The Spherepop Event Calculus}, Independent Researcher, 2024.

\bibitem{rsvpcosmology2024}
Flyxion, \textit{The RSVP Cosmological Framework: Integrating Spherepop and KES with Empirical Simulation}, Independent Researcher, 2024.

\bibitem{manifold2025}
Flyxion, \textit{The Geometry of What Is Learned: Neural Network Semantic Manifolds}, Independent Researcher, 2025.

\bibitem{synthetic2025}
Flyxion, \textit{Synthetic Grounding in Virtual Domains}, Independent Researcher, 2025.

\end{thebibliography}

\end{document}
